\section{Related Work}

Selection of auxiliary contexts has been studied mainly through the content of
the candidates. Retrieval and demonstration selection choose examples for a
frozen predictor \citep{brown2020language,min2022rethinking}, and many
selection rules score candidates by similarity to the query
\citep{liu2022what,rubin2022learning,zhang2022active,su2023selective}. Other
work studies prompt-order sensitivity and calibration
\citep{lu2022fantastically,zhao2021calibrate}, retrieval-augmented prediction
\citep{guu2020realm,lewis2020retrieval,karpukhin2020dense,borgeaud2022improving,izacard2023atlas,shi2024replug,asai2024selfrag},
and the effect of context composition on how a model uses retrieved
information \citep{liu2024lost}. Recent KBS studies likewise make external
information acquisition adaptive, including structured engineering-knowledge
retrieval \citep{siddharth2024retrieval}, causal and conflict-aware retrieval
\citep{crgs2026causal}, adaptive multi-hop evidence retrieval
\citep{acra2025retrieval}, and pipelines that adapt the granularity of the
retrieved context \citep{morenocediel2026optimising}. A second family enforces
diversity or
coverage, using penalties for repeated content or submodular objectives
\citep{carbonell1998use,gupta-etal-2023-coverage,levy2023diverse}, and listwise
rerankers compare candidate sets directly
\citep{nogueira2019passage,ma2023zero,qin2024large,sachan2022improving}. These
methods determine which candidates are relevant or complementary in content.
Relevance is computed from the candidate and the query alone, while diversity
and coverage objectives are set functions of the selected content; none of them
is defined through the effect of a candidate on the predictor's loss.

A smaller line of work values candidates through their effect on the
prediction. Outcome-aware selection estimates the incremental utility of a
demonstration relative to an empty prompt \citep{hashimoto-etal-2024-take}, and
data valuation and influence estimation attribute a model's output to training
examples
\citep{ghorbani2019data,koh2017understanding,pruthi2020estimating,ilyas2022datamodels}.
These approaches define value with respect to the empty set or to a fixed
reference, which is the standalone utility $u(j\mid\emptyset)$ in our
notation. The difficulty of using auxiliary information well is also visible in
recent KBS work on negative transfer, where task-specific information is
relearned to suppress harmful cross-task transfer \citep{ant2025negative} and
source-domain knowledge is selectively integrated before distillation
\citep{liang2026parameter}. Such methods decide what should be transferred at the task
or domain level; our object is the candidate-level marginal $m(j\mid A)$ with
respect to a nonempty, evolving selected set, and it coincides with the
standalone utility only at the first selection.

Set-selection methods supply the structural counterpart. Submodular objectives
formalize coverage, diversity, and diminishing returns
\citep{lin2011class,nemhauser1978analysis,krause2014submodular}, determinantal
point processes model diverse subsets probabilistically
\citep{kulesza2012determinantal}, and adaptive submodularity and influence
maximization extend the tools to sequential decisions
\citep{golovin2011adaptive,leskovec2007cost}. KBS has also studied adaptive
subset construction through feature-selection objectives, including
model-independent diversity-based importance \citep{mousavi2025vsi} and
graph-regularized selection under partial supervision \citep{xie2025partial};
such objectives select informative or nonredundant input dimensions, whereas
R-MUR evaluates an auxiliary context through its incremental effect on the
downstream predictor. Classical monotone formulations reward every additional
element, and a context can instead increase the loss, so the marginal utility
we estimate may be negative; this is the phenomenon that the transfer
literature calls negative transfer
\citep{pan2010survey,torrey2010transfer,wang2019characterizing}.

Learned routing treats selection as a decision problem. Mixture-of-experts
systems route inputs to expert subsets
\citep{jacobs1991adaptive,shazeer2017outrageously,zhou2022mixture}, sparse and
soft routing architectures reduce computation while increasing capacity
\citep{fedus2022switch,lewis2021base,roller2021hash,puigcerver2024sparse}, and
routing networks select computation paths across tasks
\citep{rosenbaum2018routing,ma2018modeling}. Adaptive routing has recently
appeared in KBS across several settings: expert activation conditioned on
temporal behaviour patterns \citep{chen2024bpmoe}, expert routing by attention
over multimodal state \citep{dream2026routing}, sparse expert selection for
collaborative perception driven by agent-specific information
\citep{cogmoe2026sparse}, in-context experts for sequential recommendation
\citep{sun2025adaptive}, and query-specific intervention selection from
model-behaviour signals \citep{failuremode2026routing}. These studies establish the value
of adaptive routing, but the routed object is an expert, an intervention, or a
computational component. Decision-focused learning links
prediction quality to downstream decision quality
\citep{bertsimas2020predictive,elmachtoub2022smart,wilder2019melding,mandi2024decision},
and sequential selection is related to bandit exploration and off-policy value
estimation \citep{auer2002finite,dudik2011doubly}. R-MUR belongs to this line
but routes auxiliary contexts rather than computation or model components, and
it keeps the predictor fixed so that the routing target is their
state-conditioned marginal effect on end-task loss.
